%----------------------------------------------------------------------------------------
%	PACKAGES AND OTHER DOCUMENT CONFIGURATIONS
%----------------------------------------------------------------------------------------

\documentclass[10pt,stdletter,dateno,sigleft]{newlfm} % Extra options: 'sigleft' for a left-aligned signature, 'stdletternofrom' to remove the from address, 'letterpaper' for US letter paper - consult the newlfm class manual for more options

%\usepackage{charter} % Use the Charter font for the document text
\usepackage{graphics}

\usepackage{standalone}


%% with this patch, the newlfm should work (BEGIN)
\usepackage{etoolbox}
\makeatletter
\patchcmd{\@zfancyhead}{\fancy@reset}{\f@nch@reset}{}{}
\patchcmd{\@set@em@up}{\f@ncyolh}{\f@nch@olh}{}{}
\patchcmd{\@set@em@up}{\f@ncyolh}{\f@nch@olh}{}{}
\patchcmd{\@set@em@up}{\f@ncyorh}{\f@nch@orh}{}{}
\makeatother
%% with this patch, the newlfm should work (END)

\newsavebox{\Luiuc}\sbox{\Luiuc}{\parbox[b]{1.75in}{\vspace{0.5in}
\includegraphics[width=1.2\linewidth]{/Users/hectorbahamonde/RU/Job_Market/logo.png}}} % Company/institution logo at the top left of the page
\makeletterhead{Uiuc}{\Lheader{\usebox{\Luiuc}}}

\newlfmP{sigsize=50pt} % Slightly decrease the height of the signature field
\newlfmP{addrfromphone} % Print a phone number under the sender's address
\newlfmP{addrfromemail} % Print an email address under the sender's address
\PhrPhone{\texttt{p}} % Customize the "Telephone" text
\PhrEmail{\texttt{e}} % Customize the "E-mail" text

\lthUiuc % Print the company/institution logo

% email input begin
\newread\fid
\newcommand{\readfile}[1]% #1 = filename
{\bgroup
  \endlinechar=-1
  \openin\fid=#1
  \read\fid to\filetext
  \loop\ifx\empty\filetext\relax% skip over comments
    \read\fid to\filetext
  \repeat
  \closein\fid
  \global\let\filetext=\filetext
\egroup}
\readfile{/Users/hectorbahamonde/RU/Bibliografia_PoliSci/email.txt}
% email input end

\usepackage{color} % for colors
\usepackage{hyperref} % for links

\hypersetup{
  colorlinks = true,
  urlcolor = blue,
  pdfpagemode = UseNone
}

%----------------------------------------------------------------------------------------
%	YOUR NAME AND CONTACT INFORMATION
%----------------------------------------------------------------------------------------

%\namefrom{\vspace{-3cm}Hector Bahamonde, PhD} % Name

\addrfrom{
{\vspace{-2cm}}\\ % Date
[12pt]
{\bf Hector Bahamonde, PhD}\\
{\color{blue}{\bf Assistant Professor}}\\
\emph{{\color{blue}Social Sciences Institute}}\\
{\color{blue}O$'$Higgins University}\\
Rancagua, Chile
}

\phonefrom{+1(504) 941-9131} % Phone number

\emailfrom{\href{mailto:\filetext}{\filetext} \\ 
\texttt{w}: \href{http://www.hectorbahamonde.com}{www.HectorBahamonde.com}\\
\today, download latest version \href{http://github.com/hbahamonde/Job_Market/raw/master/Bahamonde_Diversity_Statement.pdf}{\texttt{{\color{red}here}}}}% Link to last version


%----------------------------------------------------------------------------------------
%	ADDRESSEE AND GREETING/CLOSING
%----------------------------------------------------------------------------------------

%\greetto{\vspace{-1cm}Dear Members of the Search Committee,} % Greeting text
%\closeline{\vspace{-1cm}\includegraphics[width=3cm]{/Users/hectorbahamonde/RU/Administracion/signature.pdf} \hspace{-4cm}} % Closing text




%----------------------------------------------------------------------------------------

\begin{document}
\begin{newlfm}

%----------------------------------------------------------------------------------------
%	LETTER CONTENT
%----------------------------------------------------------------------------------------



\vspace{-2cm}{\bf \huge Diversity Statement}\\

I was lucky enough to attend Rutgers for grad school, one of the most diverse universities in the United States. Serving as a teaching assistant, specially in large classrooms, gave me the opportunity to experience diversity on a daily basis. 

I not only support diversity, but also enjoy being in diverse environments. Diversity makes individuals more tolerant, flexible and educated. Diversity also ``forces'' us (educators), in a good way, to develop teaching philosophies that are flexible, considerate and interesting for diverse student bodies. 

There are many types of diversity: economic, cultural, sexual, political, and more. We, as social scientists, should know that all these types are part of our daily life. Everything we do, will most likely be related to economic inequality, racial politics, or gender representation. Everything we say is subject to debate and tension. These disputes should be treated even more carefully if our audience is diverse. It is not about not making political jokes, for example. It is more about answering the following questions: \emph{How can I make today's lecture interesting enough, so my students engage not only with me but with all their classmates? How should I present this contentious issue related to race, for example, in a way that at the end of class, my students know that there is still much more to be done?} In my experience, the answers have less to do with teaching ``the facts.'' They have to do with knowing how to expose students to complex issues. \emph{Should the rich be taxed to help the poor? Was the atomic bomb a reasonable decision? Are all democratic outcomes ``good''?} (Hitler was democratically elected) \emph{Should rich states help poor states?} (economists say ``yes,'' political scientists say ``no''), etc. All these topics can be answered from so many angles. Moreover, \emph{students will answer these questions based on their own economic, cultural, sexual and political backgrounds}. At the end of the day, I know that I did my job well if I know that my students felt \emph{curious} about these issues, even at the cost of leaving the classroom in an atmosphere of ``uncertainty.''

All and all, as faculty, I think diversity always plays in our favor. In every sense, it makes lecturing more interesting. It also plays in favor of our  students. Diversity exposes them to different perspectives. Even more importantly, diversity forms better citizens better able to engage in a diverse world. That is why I think we should always ``take advantage'' of diversity.

% ending
{More information, \href{http://www.hectorbahamonde.com/teaching/}{\texttt{syllabi}}, my \href{http://github.com/hbahamonde/Job_Market/raw/master/Bahamonde_Research_Statement.pdf}{\texttt{research}}, \href{http://github.com/hbahamonde/Job_Market/raw/master/Bahamonde_Teaching_Statement.pdf}{\texttt{teaching}} and \href{http://github.com/hbahamonde/Job_Market/raw/master/Bahamonde_Diversity_Statement.pdf}{\texttt{diversity}} statements, as well as other \href{http://www.hectorbahamonde.com/research/}{\texttt{papers}} are available on my website: \href{http://www.hectorbahamonde.com}{\texttt{www.HectorBahamonde.com}}. Thank you for considering my application. I look forward to hearing from you.\unskip}

\vspace{-10cm}

\end{newlfm}
\end{document}

